\documentclass[11pt]{article}

\usepackage[letterpaper,margin=1.05in]{geometry}
\usepackage[T1]{fontenc}
\usepackage[utf8]{inputenc}
\usepackage{amsmath,amsthm,mathtools}
\usepackage{newtxtext,newtxmath}
\usepackage{bm}
\usepackage{microtype}
\usepackage{booktabs}
\usepackage{graphicx}
\usepackage{enumitem}
\usepackage[round,authoryear]{natbib}
\usepackage{xcolor}
\usepackage{hyperref}
\usepackage[nameinlink,noabbrev]{cleveref}
\usepackage{fancyhdr}
\usepackage{setspace}

\hypersetup{
  colorlinks=true,
  linkcolor=blue!45!black,
  citecolor=blue!45!black,
  urlcolor=blue!45!black,
  pdftitle={Exact Sparsity without Penalization},
  pdfauthor={Hansheng Jiang}
}
\setstretch{1.03}
\setlength{\parskip}{0.25em}
\setlength{\parindent}{1.5em}
\setlist[itemize]{leftmargin=1.4em,itemsep=0.15em,topsep=0.3em}
\setlist[enumerate]{leftmargin=1.7em,itemsep=0.15em,topsep=0.3em}
\allowdisplaybreaks[2]

\pagestyle{fancy}
\fancyhf{}
\fancyhead[L]{\small Exact sparsity by random reweighting}
\fancyhead[R]{\small \thepage}
\renewcommand{\headrulewidth}{0.3pt}
\setlength{\headheight}{14pt}

\newtheorem{theorem}{Theorem}[section]
\newtheorem{proposition}[theorem]{Proposition}
\newtheorem{lemma}[theorem]{Lemma}
\newtheorem{corollary}[theorem]{Corollary}
\newtheorem{assumption}[theorem]{Assumption}
\theoremstyle{definition}
\newtheorem{definition}[theorem]{Definition}
\theoremstyle{remark}
\newtheorem{remark}[theorem]{Remark}

\DeclareMathOperator*{\argmax}{arg\,max}
\DeclareMathOperator*{\argmin}{arg\,min}
\DeclareMathOperator{\conv}{conv}
\DeclareMathOperator{\supp}{supp}
\DeclareMathOperator{\rank}{rank}
\DeclareMathOperator{\diag}{diag}
\DeclareMathOperator{\KL}{KL}
\DeclareMathOperator{\TV}{TV}
\newcommand{\Pcal}{\mathcal P}
\newcommand{\Mcal}{\mathcal M}
\newcommand{\Ccal}{\mathcal C}
\newcommand{\Acal}{\mathcal A}
\newcommand{\Ucal}{\mathcal U}
\newcommand{\Ncal}{\mathcal N}
\newcommand{\RR}{\mathbb R}
\newcommand{\EE}{\mathbb E}
\newcommand{\PP}{\mathbb P}
\newcommand{\one}{\bm 1}
\newcommand{\norm}[1]{\lVert #1\rVert}
\newcommand{\abs}[1]{\lvert #1\rvert}
\newcommand{\ip}[2]{\langle #1,#2\rangle}
\newcommand{\vhat}{\widehat v}
\newcommand{\Ghat}{\widehat G}
\newcommand{\Gtil}{\widetilde G}
\newcommand{\vtil}{\widetilde v}
\newcommand{\eps}{\varepsilon}
\newcommand{\dd}{\,d}

\title{\bfseries Exact Sparsity without Penalization:\\
Randomly Reweighted NPMLEs for Gaussian Mixtures}
\author{Hansheng Jiang\\[-0.2em]
\small Rotman School of Management, University of Toronto}
\date{September 2026}

\begin{document}
\maketitle

\begin{abstract}
The nonparametric maximum likelihood estimator (NPMLE) for a Gaussian
location mixture optimizes over all mixing distributions.  Although this
formulation is convex in the induced density, a maximizing mixing law can be
nonunique and the classical finite-sample support bound is linear in the
sample size.  We study a randomly reweighted NPMLE obtained from a single
weighted likelihood maximization with independent
$W_i\sim\mathrm{Gamma}(\alpha_n,\alpha_n)$ weights.  Normalizing the weights
produces a Dirichlet perturbation of the empirical distribution, but we use
the perturbation as an estimator-design device rather than as a bootstrap
sample.  For a fixed compact $d$-dimensional Gaussian location-mixture model,
let
\[
 r_n=\left(\frac{\log n}{\log\log n}\right)^d,
 \qquad \bar r_n=r_n+\log n,
 \qquad \alpha_n=\bar r_n\log n.
\]
We prove that, with polynomially high probability, the reweighted NPMLE is
unique and has at most $O(\bar r_n)$ atoms.  At the same time,
$\max_i\abs{W_i-1}=O(\bar r_n^{-1/2})$, its ordinary log-likelihood is within
$O(1/\log n)$ of the unweighted optimum, and its fitted log-likelihood vector
is equally close in squared Euclidean norm.  Under correct specification, its
mixture density satisfies
\[
 H^2(f_{\Ghat_W},f_{G^*})
 =O_{\PP}\!\left(\frac{(\log n)^{d+1}}
 {n(\log\log n)^d}\right).
\]
Thus, vanishing randomization yields exact polylogarithmic sparsity without
an explicit support penalty and without sacrificing global density
estimation.  In simulations in one and two dimensions, the balanced
estimator's mean Hellinger risk is within $2.1\%$ of that of the ordinary
NPMLE across all designs, while an inverse-polynomial perturbation produces
fits that are numerically indistinguishable from the ordinary fit.  The proof
develops an effective-dimension principle for positive kernel mixtures:
low-dimensional variation of fitted values controls the support of every
extreme optimizer through Hessian sensitivity and an injective Gamma change
of variables.
\end{abstract}

\noindent\textbf{Keywords:} Bayesian bootstrap; Gaussian mixtures; Hellinger
distance; nonparametric maximum likelihood; random regularization; sparsity;
weighted likelihood.

\section{Introduction}\label{sec:introduction}

Let $X_1,\ldots,X_n\in\RR^d$ be observations from a Gaussian location
mixture.  For a mixing distribution $G$ and the standard Gaussian density
$\phi_d$, write
\[
  f_G(x)=\int \phi_d(x-\theta)\,dG(\theta).
\]
The Kiefer--Wolfowitz NPMLE \citep{kieferwolfowitz1956} maximizes
$\sum_i\log f_G(X_i)$ over a convex class of probability measures.  It is attractive precisely because it does
not specify the number of mixture components.  The same infinite-dimensional
formulation, however, leaves a structural question: how complicated must an
exact maximizing mixing law be?

Classical convex geometry guarantees the existence of a maximizer with at
most $n$ atoms, and all maximizers have the same fitted values at the sample
points; see \citet{lindsay1983geometry} and the subsequent mixture-likelihood
literature.  This deterministic $n$-point bound does not explain the much
smaller models often returned in practice.  In one dimension,
\citet{polyanskiywu2020} proved that the ordinary Gaussian-mixture NPMLE has
$O(\log n)$ atoms with high probability under sub-Gaussian sampling, a
phenomenon they called self-regularization.  In several dimensions,
quantitative polylogarithmic support control for the ordinary NPMLE remains
substantially more difficult.  Recent work of \citet{wang2026} establishes
qualitative finite support for every finite-data multivariate Gaussian NPMLE,
but does not give a polylogarithmic cardinality bound or generic uniqueness.

Beyond support size, Gaussian-mixture NPMLEs have been studied extensively as
estimators of mixing distributions, marginal densities, and empirical Bayes
rules; see, among others, \citet{laird1978}, \citet{koenkermizera2014},
\citet{zhang2009}, \citet{sahaguntuboyina2020}, and
\citet{soloffetal2025}.  The resulting density-estimation theory can be nearly
parametric up to logarithmic factors, and recent entropy characterizations
clarify the minimax benchmark \citep{jiaetal2023}.  These results establish the
statistical value of the NPMLE, but they do not provide a unique exact mixing
representation with polylogarithmic support in several dimensions.

This paper takes a different route.  Rather than forcing the equal-weight
NPMLE to reveal a sparse representation, we introduce an asymptotically
negligible random tilt of the likelihood and analyze the optimizer selected by
that tilt.  Given a concentration parameter $\alpha_n$, draw independently
\[
 W_i\sim\mathrm{Gamma}(\alpha_n,\alpha_n),
 \qquad i=1,\ldots,n,
\]
where shape and rate are both $\alpha_n$, and define
\begin{equation}\label{eq:rrnpmle-intro}
 \Ghat_W\in\argmax_{G\in\Pcal(K)}
       \sum_{i=1}^n W_i\log f_G(X_i).
\end{equation}
The weights have mean one and variance $1/\alpha_n$.  After normalization,
$(W_i/\sum_jW_j)_i$ is Dirichlet with common parameter $\alpha_n$.
Dirichlet-weighted likelihoods are classical in the Bayesian bootstrap of
\citet{rubin1981} and the weighted likelihood bootstrap of
\citet{newtonraftery1994}.  Their customary role is to generate a distribution
of bootstrap or approximate posterior draws.  Here a single, increasingly
concentrated draw is used as a point estimator.  The concentration is chosen
to expose a sparse extreme point while keeping the statistical criterion
asymptotically indistinguishable from ordinary likelihood.

The main result shows that this randomization acts as an exact regularizer.
For a fixed compact parameter set and fixed dimension, choose
\[
 \alpha_n=\left\{\left(\frac{\log n}{\log\log n}\right)^d
                 +\log n\right\}\log n.
\]
Then, with probability at least $1-Cn^{-b}$ for every fixed $b>0$, the
weighted optimizer is unique and has polylogarithmic support.  Its total
ordinary log-likelihood loss is $O(1/\log n)$, while its mixture density has
the same near-parametric Hellinger accuracy established here for all
one-unit near-maximizers.  The perturbation itself vanishes:
$\max_i\abs{W_i-1}=O_{\PP}(\bar r_n^{-1/2})$.  Moreover, for any prescribed
$L>0$, one may instead take $\alpha_n=n^{2L+2}$; the perturbation is then at
most $n^{-L}$ with overwhelming probability, yet the same exact support and
Hellinger conclusions continue to hold.

The contribution is therefore not random weighting per se.  It is the joint
estimator theorem: a single reweighted NPMLE is simultaneously
\emph{unique}, \emph{exactly sparse}, \emph{near-optimal in ordinary
likelihood}, and \emph{globally accurate in Hellinger distance}.  Randomness
resolves degeneracy, but the support theorem is not a generic tie-breaking
statement.  Its cardinality is controlled by an analytic effective dimension
of the Gaussian fitted-value set.

A numerical study makes the proximity statement operational.  Along a
regularization path, weight dispersion decays at the expected
$\alpha^{-1/2}$ scale, whereas the ordinary-likelihood loss and distance
between fitted densities decay approximately as $\alpha^{-1}$.  Across
three data-generating mechanisms and eleven sample-size--design combinations,
the balanced estimator has essentially the same Hellinger risk and support
size as the ordinary NPMLE.  By contrast, Dirichlet parameter one---the
Bayesian-bootstrap scale---does reduce support slightly in some experiments,
but incurs a substantial likelihood and risk loss.  Concentration of the
weights, rather than random weighting alone, is therefore essential.

The proof has four ingredients.  First, an extreme optimizing measure with
$k$ atoms creates a $(k-1)$-dimensional family of feasible mass perturbations.
This forces at least $k-1$ unit eigenvalues in a scaled Hessian of the weighted
value function.  Second, if all relative fitted vectors are well approximated
by an $r$-dimensional subspace, concentrated weights localize the optimizer to
that subspace.  Third, the injective map
$w_i\mapsto w_i\{1+\log(\vhat_i(w)/\vhat_i(\one))\}$ converts the Hessian
lower bound into a Jacobian determinant.  Integrating this determinant under
the Gamma density yields an exponential moment bound for exact support.
Fourth, Gaussian analyticity gives an $r=O(r_n)$ approximation by truncated
multivariate exponentials, and a generic evaluation-independence argument
turns sparse extreme points into a unique mixing law.  A finite likelihood
net then transfers the $o(1)$ ordinary likelihood gap into a global Hellinger
bound.

The rest of the paper is organized as follows.  \Cref{sec:model} defines the
estimator and states the main statistical results.  \Cref{sec:effective}
gives the abstract effective-dimension theorem.  \Cref{sec:gaussian}
specializes the geometry to Gaussian mixtures and proves uniqueness.
\Cref{sec:hellinger} develops the near-MLE Hellinger transfer.
\Cref{sec:numerics} studies finite-sample sparsity and proximity to the
ordinary NPMLE.  \Cref{sec:discussion} discusses interpretation and
limitations.  Complete proofs of the geometric, analytic, and entropy lemmas
are collected in the appendices.

\section{Model, reweighted estimator, and main results}\label{sec:model}

\subsection{Gaussian location mixtures}

Fix an integer $d\ge1$ and a nonempty compact set
$K\subset B(0,S)\subset\RR^d$, where $S<\infty$.  Let $\Pcal(K)$ denote the
Borel probability measures supported on $K$, and define
\[
 \phi_d(x)=(2\pi)^{-d/2}\exp(-\norm{x}^2/2),
 \qquad f_G(x)=\int_K\phi_d(x-\theta)\,dG(\theta).
\]
The data are generated according to
\begin{equation}\label{eq:model}
 X_i=\Theta_i+Z_i,\qquad
 \Theta_i\stackrel{\mathrm{iid}}\sim G^*,\quad
 Z_i\stackrel{\mathrm{iid}}\sim N(0,I_d),
\end{equation}
with the two sequences independent.  Correct specification means
$G^*\in\Pcal(K)$.  We use total log-likelihood
\[
 L_n(G)=\sum_{i=1}^n\log f_G(X_i)
\]
and squared Hellinger distance
\[
 H^2(p,q)=\int_{\RR^d}(\sqrt p-\sqrt q)^2\,dx.
\]
An ordinary NPMLE is any $\Ghat_0\in\argmax_{G\in\Pcal(K)}L_n(G)$.
Compactness and positivity of the Gaussian kernel guarantee existence.
Although the maximizing mixing law need not be unique, its fitted vector
$(f_{\Ghat_0}(X_i))_{i=1}^n$ is unique.

\subsection{The randomly reweighted NPMLE}

For $a>0$, let $\mathrm{Gamma}(a,a)$ denote the Gamma law with shape and
rate $a$.  Independently of the data, draw
\begin{equation}\label{eq:weights}
 W_1,\ldots,W_n\stackrel{\mathrm{iid}}\sim\mathrm{Gamma}(\alpha_n,\alpha_n)
\end{equation}
and define the reweighted NPMLE by
\begin{equation}\label{eq:rrnpmle}
 \Ghat_W\in\Mcal_X(W):=
 \argmax_{G\in\Pcal(K)}\sum_{i=1}^nW_i\log f_G(X_i).
\end{equation}
The fitted vector of all members of $\Mcal_X(W)$ is unique.  When the
maximizing measure is unique, we use $\Ghat_W$ for that measure.

Set $T_W=\sum_iW_i$, $P_i=W_i/T_W$, and
$F_n^{(\alpha_n)}=\sum_iP_i\delta_{X_i}$.  Multiplication of an objective by
the positive scalar $T_W$ does not change its maximizers, so
\[
 \Mcal_X(W)=\argmax_{G\in\Pcal(K)}
 \int\log f_G\,dF_n^{(\alpha_n)}.
\]
The next proposition records why this remains a centered perturbation of
ordinary likelihood.

\begin{proposition}[Dirichlet representation and unchanged target]
\label{prop:dirichlet}
Conditional on the data,
$(P_1,\ldots,P_n)\sim\mathrm{Dirichlet}(\alpha_n,\ldots,\alpha_n)$ and is
independent of $T_W\sim\mathrm{Gamma}(n\alpha_n,\alpha_n)$.  For every real
function $h$ on the sample points,
\[
 \EE(F_n^{(\alpha_n)}h\mid X)=F_nh,
 \qquad
 \operatorname{Var}(F_n^{(\alpha_n)}h\mid X)
 =\frac{F_nh^2-(F_nh)^2}{n\alpha_n+1},
\]
where $F_n=n^{-1}\sum_i\delta_{X_i}$.  Moreover,
\[
 \EE\bigl(\norm{F_n^{(\alpha_n)}-F_n}_{\TV}\mid X\bigr)
 \le \frac12\sqrt{\frac{n-1}{n\alpha_n+1}}.
\]
If $X_i\stackrel{\mathrm{iid}}\sim f_{G^*}$ and
$\EE_{G^*}\abs{\log f_G(X_1)}<\infty$, then
\[
 \EE_{X,W}\frac1n\sum_iW_i\log f_G(X_i)
 =\EE_{G^*}\log f_G(X_1)
 =\EE_{X,W}\sum_iP_i\log f_G(X_i).
\]
Under correct specification, the population criterion gap between $G^*$ and
$G$ is $\KL(f_{G^*}\|f_G)$.  Thus the reweighting introduces no population
support penalty and does not change the density target.
\end{proposition}

The proof is given in \Cref{app:dirichlet}.  The distinction from the
Bayesian bootstrap is quantitative and inferential.  The Bayesian bootstrap
uses common Dirichlet parameter one and treats repeated weighted optima as a
distribution.  Here $\alpha_n\to\infty$, the empirical perturbation vanishes,
and one draw defines the estimator.

\subsection{A joint sparsity and accuracy theorem}

For all sufficiently large $n$, define
\begin{equation}\label{eq:scales}
 r_n=\left(\frac{\log n}{\log\log n}\right)^d,
 \qquad \bar r_n=r_n+\log n,
 \qquad \alpha_n=\bar r_n\log n.
\end{equation}
The additional $\log n$ in $\bar r_n$ is the probability-tail cost.  It is
asymptotically absorbed by $r_n$ whenever $d\ge2$.

\begin{theorem}[Exact regularization by random reweighting]
\label{thm:main}
Assume \eqref{eq:model} with fixed $d$, fixed compact
$K\subset B(0,S)$, and $G^*\in\Pcal(K)$.  Let the weights satisfy
\eqref{eq:weights} with the concentration in \eqref{eq:scales}.  For every
$b>0$, there are constants $C_b<\infty$ and $n_0$ depending only on
$d,K,S,b$ such that, uniformly over $G^*\in\Pcal(K)$, the following event has
probability at least $1-C_bn^{-b}$ for every $n\ge n_0$:
\begin{enumerate}[label=\textup{(\roman*)}]
\item The weighted NPMLE is unique and
\[
  \#\supp(\Ghat_W)\le C_b\bar r_n.
\]
\item The perturbation is uniformly small:
\[
  \max_{1\le i\le n}\abs{W_i-1}\le \frac{C_b}{\sqrt{\bar r_n}}.
\]
\item For every ordinary NPMLE $\Ghat_0$,
\[
  0\le L_n(\Ghat_0)-L_n(\Ghat_W)\le\frac{C_b}{\log n},
\]
and
\[
 \sum_{i=1}^n
 \left\{\log\frac{f_{\Ghat_W}(X_i)}{f_{\Ghat_0}(X_i)}\right\}^2
 \le\frac{C_b}{\log n}.
\]
\item The fitted density obeys the global risk bound
\[
 H^2(f_{\Ghat_W},f_{G^*})
 \le C_b\frac{(\log n)^{d+1}}
 {n(\log\log n)^d}.
\]
\end{enumerate}
For $d\ge2$, the support bound in part \textup{(i)} is
$O\{(\log n/\log\log n)^d\}$; for $d=1$, it is $O(\log n)$.
\end{theorem}

The theorem separates a conditional structural statement from its statistical
consequence.  Parts (i)--(iii), conditional on the mild radius condition
$\max_i\norm{X_i}=O(\sqrt{\log n})$, do not require correct specification.
Part (iv) uses $G^*\in\Pcal(K)$ to compare likelihood with the truth.

The chosen concentration $\alpha_n=\bar r_n\log n$ is a convenient balanced
regime.  The next result shows that exact sparsity does not require a
perturbation of this magnitude.

\begin{corollary}[Arbitrarily small perturbations]
\label{cor:tiny}
Under the assumptions of \Cref{thm:main}, fix any $L>0$ and instead take
$\alpha_n=n^{2L+2}$.  For every $b>0$, with probability at least
$1-C_bn^{-b}$ for all sufficiently large $n$, the weighted NPMLE is unique,
\[
 \#\supp(\Ghat_W)\le C_{b,L}\bar r_n,
 \qquad
 \max_i\abs{W_i-1}\le n^{-L},
\]
and the Hellinger conclusion of \Cref{thm:main}(iv) holds.  Hence an
inverse-polynomially small perturbation can induce exact polylogarithmic
sparsity.
\end{corollary}

Hellinger convergence also identifies the mixing distribution on a compact
parameter space.

\begin{corollary}[Weak consistency of the mixing law]
\label{cor:weak}
Under \Cref{thm:main}, $\Ghat_W\Rightarrow G^*$ in probability in the weak
topology on $\Pcal(K)$.
\end{corollary}

\begin{proof}
By \Cref{thm:main}(iv), $H(f_{\Ghat_W},f_{G^*})\to0$ in probability.  Every
subsequence has a further subsequence along which this convergence holds
almost surely.  Compactness of $\Pcal(K)$ gives a further weakly convergent
subsequence, say $\Ghat_W\Rightarrow G$.  Pointwise continuity of the
Gaussian kernel implies $f_{\Ghat_W}(x)\to f_G(x)$ for every $x$, while
Hellinger convergence implies $L^1$ convergence to $f_{G^*}$.  Thus
$f_G=f_{G^*}$ almost everywhere.  Characteristic functions satisfy
$\widehat f_G(t)=e^{-\norm{t}^2/2}\widehat G(t)$; the Gaussian factor is
nonzero, so $G=G^*$.  Every subsequential limit is therefore $G^*$.
\end{proof}

\section{An effective-dimension principle for exact sparsity}
\label{sec:effective}

The support phenomenon is not specific to Gaussian kernels.  This section
states an abstract theorem for any positive continuous kernel whose fitted
vectors have small linear approximation width.  The Gaussian structure enters
only later, through a sharp bound on that width.

\subsection{Positive kernel hulls and optimizer fibers}

Let $\Theta$ be a nonempty compact metric space and let
$A:\Theta\to(0,\infty)^n$ be continuous.  Write
\[
 \Acal=A(\Theta),\qquad \Ccal=\conv(\Acal).
\]
The set $\Ccal$ is exactly the set of fitted vectors
$\int A(\theta)\,dG(\theta)$ with $G\in\Pcal(\Theta)$.  For
$w\in(0,\infty)^n$, define
\begin{equation}\label{eq:value}
 \Psi(w)=\max_{v\in\Ccal}w^\top\log v,
 \qquad
 \vhat(w)=\argmax_{v\in\Ccal}w^\top\log v,
 \qquad z(w)=\log\vhat(w),
\end{equation}
where logarithms and vector division are coordinatewise.  Strict concavity in
$v$ makes $\vhat(w)$ unique.  The corresponding optimizer fiber is
\begin{equation}\label{eq:fiber}
 \Mcal_A(w)=\left\{G\in\Pcal(\Theta):
       \int A(\theta)\,dG(\theta)=\vhat(w)\right\}.
\end{equation}
This fiber is the full set of maximizing mixing measures.

A member of a convex set is \emph{extreme} if it is not a nontrivial convex
combination of two other members.  Let
\begin{equation}\label{eq:sext}
 s_{\mathrm{ext}}(w)=
 \max\left\{\#\supp(G):G\text{ is an extreme point of }\Mcal_A(w)\right\}.
\end{equation}
The maximum is well defined and lies in $\{1,\ldots,n\}$.  Indeed, as shown in
\Cref{lem:extreme}, an optimizer is extreme exactly when it is finitely
supported and its evaluation vectors are linearly independent.

Take the equal-weight fitted vector $v_0=\vhat(\one)$ and define the relative
fitted-value set
\begin{equation}\label{eq:relative-set}
 \Ucal=\{v/v_0-\one:v\in\Ccal\}.
\end{equation}
Suppose that a rank-$r$ orthogonal projection $P$ satisfies
\begin{equation}\label{eq:width}
 \sup_{u\in\Ucal}\norm{(I-P)u}_2\le\eta.
\end{equation}
Thus $r$ is an effective linear dimension and $\eta$ is the residual width.
Both may depend on the fixed data and kernel dictionary, but not on the random
weights.

For $\tau\ge0$, call $\vtil\in\Ccal$ a \emph{$\tau$-optimal weighted fit} if
\begin{equation}\label{eq:tau-opt}
 w^\top\log\vtil\ge\Psi(w)-\tau.
\end{equation}
The tolerance is measured in total weighted log-likelihood units.

\begin{theorem}[Effective-dimension exact sparsification]
\label{thm:effective}
There are universal constants $C,c>0$ such that the following holds.  Fix
$q\ge1$, put $Q=q+\log(2n)$, and suppose
\begin{equation}\label{eq:effective-conditions}
 \alpha\ge C(r+Q),
 \qquad
 \eta\sqrt{\alpha nQ}\le1.
\end{equation}
Let $W_i\stackrel{\mathrm{iid}}\sim\mathrm{Gamma}(\alpha,\alpha)$.  With
probability at least $1-3e^{-q}$,
\begin{equation}\label{eq:effective-support}
 s_{\mathrm{ext}}(W)\le C(r+q),
\end{equation}
\begin{equation}\label{eq:effective-exact-fit}
 0\le \one^\top\log v_0-\one^\top\log\vhat(W)
 \le\frac{C(r+q)}{\alpha},
 \qquad
 \norm{\log\{\vhat(W)/v_0\}}_2^2
 \le\frac{C(r+q)}{\alpha}.
\end{equation}
On the same event, simultaneously for every $0\le\tau<c$ and every
$\tau$-optimal fit $\vtil$,
\begin{equation}\label{eq:effective-approx-fit}
 0\le \one^\top\log v_0-\one^\top\log\vtil
 \le C\left\{\frac{r+q}{\alpha}+\tau\right\},
 \qquad
 \norm{\log(\vtil/v_0)}_2^2
 \le C\left\{\frac{r+q}{\alpha}+\tau\right\}.
\end{equation}
The support conclusion concerns exact extreme optimizers.  No support bound is
asserted for arbitrary approximate fitted measures.
\end{theorem}

The theorem reveals the scale of the randomization.  The shape parameter need
only dominate the effective dimension and the logarithmic tail cost.  Larger
$\alpha$ makes the objective closer to equal weighting, but the residual
width must then be approximated more accurately.  Gaussian analyticity allows
both conditions to hold even for polynomially large $\alpha$.

\subsection{Why support is visible to the Hessian}

We summarize the geometric mechanism; full proofs appear in
\Cref{app:effective}.  The value function $\Psi$ is convex and continuously
differentiable with $\nabla\Psi=z$, and $z$ is locally Lipschitz.  Hence the
Hessian exists almost everywhere.  If an extreme optimizer has $k$ atoms,
then changing its $k$ masses while preserving their sum creates a
$(k-1)$-dimensional space of two-sided feasible perturbations of the fitted
vector.  A second-order comparison gives the following matrix inequality.

\begin{lemma}[Support--sensitivity inequality]
\label{lem:hessian-summary}
For almost every $w\in(0,\infty)^n$, there is an orthogonal projection
$P_*(w)$ of rank $s_{\mathrm{ext}}(w)-1$ such that
\begin{equation}\label{eq:hessian-summary}
 \diag(w)^{1/2}\nabla^2\Psi(w)\diag(w)^{1/2}\succeq P_*(w).
\end{equation}
The exceptional set depends only on the dictionary and not on a selected
optimizer.
\end{lemma}

The matrix inequality is stronger than a trace bound.  It yields a determinant
factor exponential in support size.  Let
\[
 t(w)=z(w)-z(\one),
 \qquad Z(w)=(w-\one)^\top t(w),
\]
and consider the locally Lipschitz map
\begin{equation}\label{eq:change-map}
 \mathcal F_i(w)=w_i\{1+t_i(w)\}.
\end{equation}
On the event where the weighted fit is close to $v_0$, this map is injective.
The determinant of its derivative contains a factor
$(17/9)^{s_{\mathrm{ext}}(w)-1}$.  Comparing the Gamma density at
$\mathcal F(w)$ and $w$, then applying the injective area formula
\citep{evansgariepy2015}, gives
\begin{equation}\label{eq:moment-summary}
 \EE\left[\bm 1_{\{W\in E\}}
 \left(\frac{17}{9}\right)^{s_{\mathrm{ext}}(W)-1}
 e^{-\alpha\{Z(W)+\norm{t(W)}_2^2\}}\right]\le1
\end{equation}
for a suitable localization event $E$.  The width condition
\eqref{eq:width} bounds the exponential cost in \eqref{eq:moment-summary} by
$O(r+q)$, and Markov's inequality then proves
\eqref{eq:effective-support}.  This change-of-variables step is the source of
\emph{exact} sparsity: it controls cardinality itself, rather than the size of
small mixture weights or an approximation error.

\section{Gaussian effective dimension and generic uniqueness}
\label{sec:gaussian}

We now specialize \Cref{thm:effective} to the Gaussian evaluation dictionary
\[
 A(\theta)=\bigl(\phi_d(x_1-\theta),\ldots,
                    \phi_d(x_n-\theta)\bigr),
 \qquad \theta\in K.
\]
Throughout this section, the dataset $x=(x_1,\ldots,x_n)$ is fixed and
$T=\max_i\norm{x_i}$.

\subsection{Analytic low-rank approximation}

The Gaussian identity
\[
 \phi_d(x_i-\theta)=\phi_d(x_i)
 \exp\{x_i^\top\theta-\norm{\theta}^2/2\}
\]
reduces the fitted-value geometry to multivariate exponential functions.
Truncating their Taylor series gives the required effective dimension.

\begin{proposition}[Gaussian relative-width bound]
\label{prop:gaussian-width}
Let $K\subset B(0,S)$ and let $G_0\in\Pcal(K)$ be arbitrary.  Put
$v_{0,i}=f_{G_0}(x_i)$.  For every integer $m\ge0$, there is an orthogonal
projection $P_m$ of rank at most
\begin{equation}\label{eq:rank-m}
 1+\binom{m+d}{d}
\end{equation}
whose range contains $\one$ and such that
\begin{equation}\label{eq:eta-m}
 \sup_{G\in\Pcal(K)}
 \norm{(I-P_m)\{f_G(x)/v_0-\one\}}_2
 \le \eta_m:=\sqrt n\,\eps_m,
\end{equation}
where
\begin{equation}\label{eq:eps-m}
 \eps_m=\exp(2ST+S^2/2)\frac{(ST)^{m+1}}{(m+1)!}.
\end{equation}
When $ST=0$, one may take $\eps_m=0$.
\end{proposition}

If $T\le C_0\sqrt{\log n}$ and
$m\asymp\log n/\log\log n$, then the rank in \eqref{eq:rank-m} is
$O(r_n)$, while $\eps_m$ is smaller than any prescribed inverse power of
$n$ after increasing the constant in $m$.  Thus the width condition in
\Cref{thm:effective} holds for a wide range of $\alpha$.

\begin{corollary}[Conditional structural and likelihood bounds]
\label{cor:conditional}
Fix $d,S,C_0<\infty$, a nonempty compact $K\subset B(0,S)$, and a dataset
with $T\le C_0\sqrt{\log n}$.  Let $r_n,\bar r_n$ be as in
\eqref{eq:scales}, set $\alpha_n=\bar r_n\log n$, and draw independent
$\mathrm{Gamma}(\alpha_n,\alpha_n)$ weights.  For every $b>0$, with
conditional probability at least $1-C_bn^{-b}$,
\[
 s_{\mathrm{ext}}(W)\le C_b\bar r_n,
 \qquad
 \max_i\abs{W_i-1}\le C_b\bar r_n^{-1/2},
\]
and every exact weighted maximizer satisfies the two bounds in
\Cref{thm:main}(iii).  These conclusions are deterministic in the data and
make no sampling or correct-specification assumption.
\end{corollary}

\subsection{From sparse extreme points to a unique measure}

A support bound on extreme points does not by itself imply that every
optimizer is sparse: a nontrivial convex combination of extreme optimizers
can have a larger or even infinite support.  For Gaussian evaluations,
generic data rule out this possibility once all extreme optimizers are
sufficiently small.

\begin{proposition}[Generic evaluation independence]
\label{prop:generic}
For each $n,d$, there is a Borel set $\mathcal X_{n,d}\subset(\RR^d)^n$ whose
complement has Lebesgue measure zero such that the following holds
simultaneously for every $x\in\mathcal X_{n,d}$.  If
$k(d+1)\le n$ and $\theta_1,\ldots,\theta_k\in\RR^d$ are distinct, then the
$n\times k$ matrix
\begin{equation}\label{eq:evaluation-matrix}
 \bigl(\phi_d(x_i-\theta_j)\bigr)_{i\le n,j\le k}
\end{equation}
has linearly independent columns.  The locations may depend on the data.
Consequently, if every extreme point of a Gaussian optimizer fiber has at
most $m$ atoms and $2m(d+1)\le n$, then that fiber is a singleton and its
unique member has at most $m$ atoms.
\end{proposition}

The key point is simultaneity over the uncountable collection of possible
locations.  The proof treats a linear dependence as an analytic incidence
set with $k(d+1)-1$ nuisance parameters.  Since this dimension is below the
number $n$ of sample equations, its projection onto the $nd$ data coordinates
has Lebesgue measure zero.  The uniqueness conclusion follows by comparing
two extreme optimizers and applying Krein--Milman; details are in
\Cref{app:generic}.

\begin{proof}[Proof of the structural parts of \Cref{thm:main}]
Under \eqref{eq:model}, the joint distribution of
$(X_1,\ldots,X_n)$ has a density, so the generic event in
\Cref{prop:generic} has probability one.  A Gaussian union bound gives
$\max_i\norm{X_i}\le C_b\sqrt{\log n}$ outside an event of probability
$O(n^{-b-2})$, uniformly over $G^*\in\Pcal(K)$.  Conditional on this radius
event, \Cref{cor:conditional} bounds every extreme optimizer by
$C_b\bar r_n$ atoms.  Since $\bar r_n=o(n)$, eventually
$2C_b\bar r_n(d+1)\le n$.  The second part of \Cref{prop:generic} then makes
the weighted optimizer fiber a singleton, proving \Cref{thm:main}(i).
The weight, likelihood, and fitted-vector bounds are exactly those of
\Cref{cor:conditional}.  Integrating the conditional probability bound and
combining the exceptional events proves parts (ii)--(iii).
\end{proof}

\begin{proof}[Proof of \Cref{cor:tiny} except for Hellinger accuracy]
Take $\alpha_n=n^{2L+2}$.  In \Cref{prop:gaussian-width}, choose
$m=\lceil C_L\log n/\log\log n\rceil$ with $C_L$ sufficiently large.  On
the same data-radius event as above,
$\rank(P_m)=O_L(r_n)$ and $\eps_m\le n^{-L-3}$.  With
$q=(b+2)\log n$, the two conditions in \eqref{eq:effective-conditions} hold
for all large $n$.  Hence every extreme optimizer has
$O_{b,L}(\bar r_n)$ atoms.  Generic evaluation independence again gives
uniqueness.  The Gamma concentration inequality
\[
 \PP\{\max_i\abs{W_i-1}>t\}\le2n\exp(-\alpha_nt^2/4),
 \qquad 0<t\le1,
\]
with $t=n^{-L}$ yields a bound $2ne^{-n^2/4}$.
\end{proof}

\section{Hellinger accuracy from near-maximal likelihood}
\label{sec:hellinger}

The geometric theorem gives a total ordinary likelihood gap of order
$1/\log n$.  This section converts that deterministic near-optimality into a
global density guarantee.  The argument is deliberately uniform: one data
event controls every candidate whose total log-likelihood is within one unit
of the truth.

\begin{theorem}[Uniform Hellinger control for near-MLEs]
\label{thm:near-mle}
Fix $d\ge1$ and a nonempty compact $K\subset B(0,S)$.  For every $b>0$,
there are constants $C_b,n_0$ such that, uniformly over
$G^*\in\Pcal(K)$, with probability at least $1-C_bn^{-b}$ the following
holds simultaneously for all $G\in\Pcal(K)$:
\begin{equation}\label{eq:near-truth}
 L_n(G)\ge L_n(G^*)-1
 \quad\Longrightarrow\quad
 H^2(f_G,f_{G^*})
 \le C_b\frac{(\log n)^{d+1}}
 {n(\log\log n)^d}.
\end{equation}
\end{theorem}

The proof uses a deterministic finite net $\Ncal_n$ of the compact Gaussian
mixture class.  It simultaneously approximates Hellinger distance globally
and log densities on the sample-radius ball, with
\begin{equation}\label{eq:net-entropy-summary}
 \log\abs{\Ncal_n}
 \le C\frac{(\log n)^{d+1}}{(\log\log n)^d}.
\end{equation}
For a fixed net density $p$, the square-root likelihood ratio satisfies
\[
 \EE_{G^*}\exp\left\{\frac12\sum_{i=1}^n
          \log\frac{p(X_i)}{f_{G^*}(X_i)}\right\}
 =\left(1-\frac12H^2(p,f_{G^*})\right)^n.
\]
A union bound over the net therefore excludes every Hellinger-distant
candidate with likelihood ratio above $-2$.  The complete construction and
proof appear in \Cref{app:hellinger}.

\begin{proof}[Completion of the proof of \Cref{thm:main}]
On the event of \Cref{thm:main}(iii), for all sufficiently large $n$,
\[
 L_n(\Ghat_W)\ge L_n(\Ghat_0)-1\ge L_n(G^*)-1,
\]
because $G^*$ is feasible for the ordinary likelihood.  Intersect this event
with the common data event in \Cref{thm:near-mle}.  The implication
\eqref{eq:near-truth} then applies to the weighted optimizer and proves part
(iv).  A union bound preserves probability $1-C_bn^{-b}$ uniformly over
$G^*\in\Pcal(K)$.
\end{proof}

\begin{proof}[Completion of the proof of \Cref{cor:tiny}]
The effective-dimension theorem also gives an ordinary total likelihood gap
$C_{b,L}\bar r_n/n^{2L+2}=o(1)$.  Hence the weighted optimizer satisfies
\eqref{eq:near-truth} for all sufficiently large $n$.  Apply
\Cref{thm:near-mle} and combine the exceptional events.
\end{proof}

The approximate-fit part of \Cref{thm:effective} has a direct statistical
interpretation.

\begin{corollary}[Robustness to numerical optimization]
\label{cor:approx}
Under the setting and weighting of \Cref{thm:main}, suppose a fitted vector
$\vtil$ is attainable by a mixing law $\Gtil\in\Pcal(K)$ and has total
weighted log-likelihood gap $\tau_n\le c_0/\log n$, where $c_0$ is a fixed
sufficiently small constant.  Then, with probability at least $1-C_bn^{-b}$,
\[
 0\le L_n(\Ghat_0)-L_n(\Gtil)\le C_b/\log n,
 \qquad
 \sum_i\left\{\log\frac{f_{\Gtil}(X_i)}
 {f_{\Ghat_0}(X_i)}\right\}^2\le C_b/\log n,
\]
and $\Gtil$ obeys the Hellinger bound in \Cref{thm:main}(iv).  This result
does not assert exact sparsity of an arbitrary approximate optimizer.
\end{corollary}

\begin{proof}
Apply \eqref{eq:effective-approx-fit} with the Gaussian width construction
used in \Cref{cor:conditional}.  The ordinary likelihood gap is eventually
below one, so \Cref{thm:near-mle} applies.
\end{proof}


\section{Numerical evidence: sparsity and proximity to the ordinary NPMLE}
\label{sec:numerics}

The theory is asymptotic and its principal conclusion is structural.  The
purpose of this section is therefore not to establish a new empirical risk
advantage over the ordinary NPMLE.  Instead, we examine the finite-sample
consequences that make the reweighted estimator useful: whether it retains the
ordinary estimator's density accuracy and support complexity, how quickly it
approaches the ordinary fit as the reweighting vanishes, and whether the
concentration of the weights is essential.  All comparisons are paired on the
same simulated dataset.

Numerical means, standard errors, percentages, quantiles, and regression
coefficients are reported at the displayed precision.  Range endpoints
described below as rounded extrema are the rounded minimum and maximum cell
summaries, not literal enclosing bounds on the unrounded values.

\subsection{Design and certified computation}

We use three correctly specified Gaussian location-mixture models.  The first
is the one-dimensional discrete mixture
\[
 G^*=0.25\delta_{-2.5}+0.50\delta_0+0.25\delta_{2.5},
 \qquad K=[-4,4].
\]
The second has $G^*=\operatorname{Unif}[-2.5,2.5]$ on the same parameter
space; this design checks that the conclusions are not tied to finite true
support.  The third is two-dimensional, with equal mass at the four corners
$(\pm1.75,\pm1.75)$ and $K=[-3.5,3.5]^2$.  We use
$n\in\{200,500,1000,2000\}$ and 100 replications for each one-dimensional
design, and $n\in\{250,500,1000\}$ and 40 replications for the
two-dimensional design.

Four estimators are compared.  ``Ordinary'' is the equal-weight NPMLE.
``Balanced'' uses the concentration in \eqref{eq:scales},
``Tiny'' uses $\alpha_n=n^3$, corresponding to \Cref{cor:tiny} with
$L=1/2$, and ``BB-scale'' uses $\alpha=1$ as a deliberately noisy control.
The last procedure is not proposed as a point estimator; it isolates what
happens when random weighting is used without increasing concentration.  We
normalize every realization to
\[
 \widetilde W_i=\frac{W_i}{n^{-1}\sum_jW_j}=nP_i,
\]
which leaves the optimizer unchanged and removes irrelevant common scaling.

We report numerical support size, squared Hellinger risk
$H^2(f_{\widehat G},f_{G^*})$, squared Hellinger distance to the ordinary fit,
total ordinary-likelihood loss
$L_n(\widehat G_0)-L_n(\widehat G_W)$, and
\[
 \sum_{i=1}^n\left\{
 \log f_{\widehat G_W}(X_i)-\log f_{\widehat G_0}(X_i)
 \right\}^2.
\]
Hellinger integrals are computed on $[-10,10]$ using 6001 points in one
dimension and on $[-8,8]^2$ using a $141\times141$ grid in two dimensions.
For descriptive support counts, atoms separated by less than $0.01$ in one
dimension or $0.05$ in two dimensions are treated as numerically unresolved.
This convention affects only reporting, not optimization.

The estimator is computed by fully corrective column generation.  Given a
current mixing law $G$, the global optimality residual is
\begin{equation}\label{eq:numerical-kkt}
 \mathcal R_w(G)=
 \sup_{\theta\in K}\left[
 \frac{1}{\sum_iw_i}\sum_{i=1}^n
 w_i\frac{\phi_d(X_i-\theta)}{f_G(X_i)}-1
 \right].
\end{equation}
At an exact optimum this quantity is nonpositive.  We globally search the
one-dimensional oracle on a dense mesh followed by local refinement; in two
dimensions we use a rectangular mesh and multistart continuous refinement.
After column generation, we jointly polish all atom locations and masses and
then reevaluate \eqref{eq:numerical-kkt}.  Across the 4160 fitted optimization
problems reported below, the largest recorded residual was at most
$1.86\times10^{-6}$, a conservatively upward-rounded bound; the
one-dimensional maximum rounded to $4.84\times10^{-7}$.  As a separate
optimization check, ten dispersed random
initializations on representative one- and two-dimensional datasets converged
to the same resolved support.  The largest squared Hellinger discrepancy from
the reference solution was respectively $7.9\times10^{-16}$ and
$3.7\times10^{-13}$.  These diagnostics validate the numerical comparison;
they are not used as a substitute for the theoretical uniqueness result.

\subsection{A regularization path}

We first hold fixed one sample of size $n=500$ from the three-point mixture and
vary $\alpha$ from 1 to $n^3$.  At each value we use 30 independent weight
draws.  \Cref{fig:path} reports medians and 10--90\% intervals.  The balanced
value is $\alpha_n=59.76$ for this experiment.

\begin{figure}[t]
\centering
\begin{minipage}[t]{0.48\textwidth}
  \centering\includegraphics[width=\linewidth]{figures/path_weight_deviation.pdf}
  \par\vspace{-0.2em}\centerline{\small (a) Weight perturbation}
\end{minipage}\hfill
\begin{minipage}[t]{0.48\textwidth}
  \centering\includegraphics[width=\linewidth]{figures/path_hellinger_to_ordinary.pdf}
  \par\vspace{-0.2em}\centerline{\small (b) Distance to the ordinary fit}
\end{minipage}

\par\medskip
\begin{minipage}[t]{0.48\textwidth}
  \centering\includegraphics[width=\linewidth]{figures/path_likelihood_gap.pdf}
  \par\vspace{-0.2em}\centerline{\small (c) Ordinary-likelihood loss}
\end{minipage}\hfill
\begin{minipage}[t]{0.48\textwidth}
  \centering\includegraphics[width=\linewidth]{figures/path_support.pdf}
  \par\vspace{-0.2em}\centerline{\small (d) Numerical support}
\end{minipage}
\caption{Regularization path on a fixed $n=500$ sample from the three-point
mixture.  Curves show medians across 30 independent Gamma-weight draws and
bands show the 10th and 90th percentiles.  The dashed and dotted vertical
lines mark the balanced concentration and $n^3$, respectively.}
\label{fig:path}
\end{figure}

The path displays two distinct scales.  A log--log regression of the medians
for $\alpha\ge10$ gives slope $-0.503$ for
$\max_i|\widetilde W_i-1|$, matching the $\alpha^{-1/2}$ dispersion of Gamma
weights.  The slopes are $-1.016$ for squared Hellinger distance to the
ordinary fit, $-1.016$ for ordinary-likelihood loss, and $-1.012$ for the
squared fitted-log discrepancy.  Thus the fitted vector moves at first order
in the weight perturbation while likelihood and squared density discrepancies
are locally second order.

At the balanced concentration, the median squared Hellinger distance to the
ordinary fit is $4.40\times10^{-5}$ and the median total ordinary-likelihood
loss is $0.0453$.  The ordinary fitted density has squared Hellinger risk
$0.002662$ in this sample, compared with median risk $0.002664$ for the
balanced estimator.  At $\alpha=n^3$, the corresponding distance and
likelihood loss fall to $1.52\times10^{-11}$ and $1.65\times10^{-8}$, and the
risk agrees with the ordinary risk to the displayed precision.  The median
resolved support is three atoms at both concentrations.  At the balanced
value, 28 of the 30 draws have three atoms and two have four; every tiny-regime
draw has three atoms.  By comparison, $\alpha=1$ yields median distance
$1.95\times10^{-3}$ and median likelihood loss $2.03$, with a 90th-percentile
support of four atoms.

\begin{figure}[t]
\centering
\begin{minipage}[t]{0.49\textwidth}
  \centering\includegraphics[width=\linewidth]{figures/representative_density.pdf}
  \par\vspace{-0.2em}\centerline{\small (a) Fitted densities}
\end{minipage}\hfill
\begin{minipage}[t]{0.49\textwidth}
  \centering\includegraphics[width=\linewidth]{figures/representative_support.pdf}
  \par\vspace{-0.2em}\centerline{\small (b) Mixing measures}
\end{minipage}
\caption{Representative fits from the regularization-path experiment.  The
ordinary, balanced, and tiny fitted densities are visually indistinguishable,
although their atom locations and masses need not coincide exactly.  The
BB-scale fit exhibits a visible displacement.}
\label{fig:representative}
\end{figure}

\Cref{fig:representative} also illustrates why proximity should be judged
primarily through the induced density, fitted vector, and likelihood.  At the
balanced concentration, the median one-dimensional Wasserstein distance
between the mixing laws is $0.0308$, visibly larger than the discrepancy
between the convolved densities.  Gaussian deconvolution is an ill-posed
inverse problem, so a small density perturbation need not imply comparable
finite-sample atom-by-atom stability.  This is consistent with the quantities
controlled by \Cref{thm:main}.

\subsection{Monte Carlo risk, support, and likelihood}

\Cref{tab:numerical-summary} reports the largest sample size in each design;
entries are means with Monte Carlo standard errors in parentheses.  The full
sample-size trajectories are shown in \Cref{fig:risk}.

\begin{table}[t]
\centering
\caption{Finite-sample comparison with the ordinary NPMLE.  ``Balanced'' uses
\eqref{eq:scales}, ``Tiny'' uses $\alpha_n=n^3$, and ``BB-scale'' uses
$\alpha=1$.  Entries are Monte Carlo means with standard errors in
parentheses.  Here $\Delta L_n=L_n(\widehat G_0)-L_n(\widehat G_W)$.}
\label{tab:numerical-summary}
\footnotesize
\setlength{\tabcolsep}{3.2pt}
\input{results/numerical_summary.tex}
\end{table}

\begin{figure}[t]
\centering
\begin{minipage}[t]{0.32\textwidth}
  \centering\includegraphics[width=\linewidth]{figures/risk_three_point.pdf}
  \par\vspace{-0.2em}\centerline{\small (a) Three-point, $d=1$}
\end{minipage}\hfill
\begin{minipage}[t]{0.32\textwidth}
  \centering\includegraphics[width=\linewidth]{figures/risk_uniform.pdf}
  \par\vspace{-0.2em}\centerline{\small (b) Uniform mixing, $d=1$}
\end{minipage}\hfill
\begin{minipage}[t]{0.32\textwidth}
  \centering\includegraphics[width=\linewidth]{figures/risk_square_2d.pdf}
  \par\vspace{-0.2em}\centerline{\small (c) Four-point square, $d=2$}
\end{minipage}
\caption{Mean squared Hellinger risk as a function of sample size.  Shaded
bands are pointwise 95\% Monte Carlo intervals.}
\label{fig:risk}
\end{figure}

The balanced estimator tracks the ordinary NPMLE closely in every design.
Across all eleven design--sample-size cells, the ratio of mean Hellinger risks
has rounded minimum $1.001$ and rounded maximum $1.020$; its average across
cells is $1.013$.  The mean squared Hellinger distance between the two fitted
densities has rounded minimum $8.31\times10^{-6}$ and rounded maximum
$1.76\times10^{-4}$.  The balanced and ordinary
resolved support counts agree in $90.3\%$ of the 920 paired replications, and
the mean support difference in every cell lies between $-0.125$ and $0.025$
atoms.  Thus the reweighting does not create a larger fitted model in finite
samples, but neither should it be interpreted as a device whose purpose is to
use fewer atoms than an ordinary NPMLE that happens already to be sparse.
Its contribution is to provide a generic unique sparse selection with a
quantitative guarantee.

The ordinary-likelihood comparison agrees with the scale of
\Cref{thm:main}(iii).  The balanced mean total loss has rounded minimum
$0.0345$ and rounded maximum $0.0964$.  Multiplying the cell means by $\log n$
gives rounded minimum $0.253$ and rounded maximum $0.532$; the corresponding
values for the squared fitted-log discrepancy have rounded minimum $0.479$
and rounded maximum $0.928$.  These constants are stable across sample sizes
and dimensions over the range studied.

The tiny regime gives a sharper finite-sample approximation.  Its ratio of
mean Hellinger risks to ordinary risk has rounded minimum $0.999985$ and
rounded maximum $1.000018$, and its mean squared Hellinger distance to the
ordinary fit has rounded minimum $1.43\times10^{-13}$ and rounded maximum
$9.60\times10^{-10}$.  It has exactly the same
resolved support count as the ordinary estimator in every one of the 920
paired replications.  This is direct numerical evidence for the central
interpretation of \Cref{cor:tiny}: exact sparse selection can persist while
the objective and fitted density become essentially indistinguishable from
their equal-weight counterparts.

The BB-scale control behaves differently.  Its ratio of mean Hellinger risk
to ordinary risk has rounded minimum $1.60$ and rounded maximum $2.19$, and
its mean total ordinary-likelihood loss has rounded minimum $2.14$ and
rounded maximum $7.86$.  It sometimes uses slightly
fewer atoms, but this reduction is accompanied by a material statistical
cost.  Randomization alone is therefore not the relevant regularizer; the
key is a continuous perturbation whose concentration increases sufficiently
quickly with $n$.

Taken together, the experiments support three conclusions.  First, the
balanced estimator retains the finite-sample density accuracy and support
complexity of the ordinary NPMLE.  Second, the inverse-polynomial regime can
make the two estimators operationally indistinguishable while retaining the
theoretical exact-sparsity conclusion.  Third, the proximity is naturally a
statement about fitted densities and likelihoods, not exact agreement of atom
locations in a finite-sample deconvolution problem.

\section{Discussion}\label{sec:discussion}

\paragraph{The main discovery.}
The central result is best viewed as an estimator theorem rather than a
perturbation lemma.  Random reweighting constructs a new point estimator of
the mixing law with four simultaneous properties: exact polylogarithmic
support, generic uniqueness, asymptotically negligible ordinary likelihood
loss, and global Hellinger convergence.  Each part matters.  Sparsity alone
could select a poor density; Hellinger consistency alone would not resolve the
nonuniqueness and computational burden of the mixing representation.  Their
combination is what makes the reweighted estimator statistically meaningful.

\paragraph{Regularization without a penalty.}
No term depending on $\#\supp(G)$ appears in either the sample objective or
the population target.  The regularization is geometric.  A small continuous
random tilt makes high-support optimizer fibers occupy little Gamma volume,
and the analytic effective dimension of Gaussian evaluations determines the
remaining cardinality.  In this sense, randomization behaves like an exact
support regularizer without introducing an explicit regularization bias.

\paragraph{Auxiliary randomness.}
The guarantee requires only one draw of the weights; repeated bootstrap
optimization is neither part of the definition nor needed for the theory.
The auxiliary seed can therefore be fixed and recorded.  Different successful
draws may select different sparse mixing laws, but every such draw is close to
the ordinary fitted vector and lies in the same shrinking Hellinger
neighborhood of the truth.  The randomization is thus a reproducible
regularization device rather than an uncertainty-quantification procedure.

\paragraph{Relation to self-regularization.}
The one-dimensional ordinary NPMLE is already known to have logarithmic
support under broad sampling conditions \citep{polyanskiywu2020}.  Our result
does not prove the analogous statement for the equal-weight multivariate
NPMLE\@.  Instead, it shows that an arbitrarily close objective possesses a
unique optimizer with a quantitative polylogarithmic support bound in every
fixed dimension.  The distinction is substantive: the upper support bounds
proved here do not pass to the equal-weight limit by lower semicontinuity.

\paragraph{Statistical rate.}
The Hellinger bound is near parametric up to logarithmic factors, but it is not
claimed to be sharp.  For compact or sub-Gaussian Gaussian mixtures, known
lower bounds are of order $(\log n)^d/n$ under squared Hellinger loss
\citep{kimguntuboyina2022}, and local-entropy results sharpen the broader
minimax picture \citep{jiaetal2023}.  Improving the finite-net argument while
preserving a simultaneous guarantee for all near-MLE candidates is an
important direction.

\paragraph{Computation.}
The theory is structural, not an algorithmic convergence theorem.  It implies
that an exact extreme optimum has only polylogarithmically many atoms and that
small objective error preserves fitted-value and Hellinger accuracy.  The
numerical study shows that a fully corrective column-generation method,
followed by joint location--mass polishing, can attain small global
Kiefer--Wolfowitz residuals while maintaining the predicted sparse
representation.  This supports exchange, vertex-direction, and adaptive
discretization methods, but does not establish a general complexity bound for
them.  Exploiting the Hessian response directions to obtain such a bound
remains an important algorithmic question.

\paragraph{Extensions.}
The effective-dimension theorem applies to general positive kernels.  To
extend the full result, one needs (i) a low-rank approximation of relative
fitted values, (ii) generic evaluation independence or another uniqueness
mechanism, and (iii) a likelihood-to-risk transfer for the induced density
class.  Exponential families and heteroscedastic Gaussian mixtures are natural
candidates.  Derandomization, data-adaptive choice of $\alpha_n$, and the use
of multiple sparse draws for uncertainty quantification are also open.

\appendix

\section{Proof of the effective-dimension theorem}
\label{app:effective}

\subsection{Basic convex geometry}

\begin{lemma}[Regularity of the fitted-vector problem]
\label{lem:regularity}
In the setup of \eqref{eq:value}, $\Ccal$ is compact,
$\Psi$ is convex and continuously differentiable on $(0,\infty)^n$, and
\[
 \nabla\Psi(w)=z(w)=\log\vhat(w).
\]
The maps $w\mapsto\vhat(w)$ and $w\mapsto z(w)$ are locally Lipschitz.
All optimizing measures have fitted vector $\vhat(w)$.
\end{lemma}

\begin{proof}
Since $\Acal$ is compact in $\RR^n$, Carath\'eodory's theorem expresses
every point of $\conv(\Acal)$ as a convex combination of at most $n+1$
points of $\Acal$.  Hence $\Ccal$ is compact.  Positivity and compactness give
constants $0<c_0\le C_0<\infty$ such that
$c_0\le v_i\le C_0$ for all $v\in\Ccal$ and every $i$.  The objective
$v\mapsto\sum_iw_i\log v_i$ is continuous and strictly concave, so it has a
unique maximizing fitted vector.

Fix a compact box of weights in the positive orthant and let $a_0>0$ be a
common lower bound on its coordinates.  Put
$v=\vhat(w)$ and $v'=\vhat(w')$.  Adding the two variational inequalities at
$v$ and $v'$ gives
\[
 0\le\{\nabla f_w(v)-\nabla f_{w'}(v')\}^\top(v-v'),
 \qquad f_w(u)=\sum_iw_i\log u_i.
\]
Strong concavity of $f_w$ on $[c_0,C_0]^n$ and
$\nabla f_w(v')-\nabla f_{w'}(v')=(w-w')/v'$ imply
\[
 0\le-\frac{a_0}{C_0^2}\norm{v-v'}_2^2
       +\frac1{c_0}\norm{w-w'}_2\norm{v-v'}_2.
\]
Thus $\vhat$ is locally Lipschitz, and so is its coordinatewise logarithm.

The maximum of linear functions of $w$ is convex.  For small $h$, optimality
at $w$ and $w+h$ gives
\[
 h^\top z(w)\le\Psi(w+h)-\Psi(w)\le h^\top z(w+h).
\]
Local continuity of $z$ yields differentiability of $\Psi$ with gradient
$z$.  Finally, optimizing over measures depends on a measure only through its
attainable fitted vector in $\Ccal$, proving the last assertion.
\end{proof}

\begin{lemma}[Extreme optimizer characterization]
\label{lem:extreme}
For every $w>0$, the fiber $\Mcal_A(w)$ in \eqref{eq:fiber} is nonempty,
compact, and convex.  A measure $G\in\Mcal_A(w)$ is extreme if and only if
\[
 G=\sum_{j=1}^kp_j\delta_{\theta_j},\qquad p_j>0,
\]
for finitely many distinct points whose evaluation vectors
$A(\theta_1),\ldots,A(\theta_k)$ are linearly independent.  Consequently,
every extreme optimizer has at most $n$ atoms.  The function
$s_{\mathrm{ext}}$ in \eqref{eq:sext} is Borel measurable.
\end{lemma}

\begin{proof}
Probability measures on compact $\Theta$ form a compact convex set under weak
convergence, and the continuous moment constraints defining
$\Mcal_A(w)$ are closed.  Nonemptiness follows from the representation of
$\vhat(w)\in\Ccal$ by a finite convex combination.

Let $\widehat v=\vhat(w)$ and define
\[
 b_i=\frac{w_i}{\widehat v_i\sum_{\ell=1}^nw_\ell},
 \qquad q(\theta)=b^\top A(\theta).
\]
The directional derivative from $\widehat v$ toward $A(\theta)$ gives
$q(\theta)\le1$.  If $G$ is in the fiber, then
$\int q\,dG=b^\top\widehat v=1$; hence $q=1$ on $\supp(G)$.  Therefore, on
an optimizing support, linear dependence and affine dependence of the
evaluation vectors are equivalent: applying $b^\top$ to a linear relation
shows that its coefficients sum to zero.

Suppose first that $G$ has at least $n+1$ support points, including the case of
infinite support.  Choose $n+1$ disjoint Borel neighborhoods
$E_1,\ldots,E_{n+1}$ of positive $G$-mass.  The vectors
$v_j=\int_{E_j}A\,dG\in\RR^n$ are linearly dependent.  For nonzero
coefficients $c_j$ with $\sum_jc_jv_j=0$, define
$dH=\sum_jc_j\bm 1_{E_j}\,dG$.  Then $\int A\,dH=0$ and, because $q=1$
$G$-almost surely,
$H(\Theta)=\int q\,dH=0$.  For sufficiently small $t>0$, both $G+tH$ and
$G-tH$ are distinct probability measures in the same fiber, so $G$ is not
extreme.

If $G$ has finite support but the evaluation vectors are dependent, take a
nonzero relation $\sum_jc_jA(\theta_j)=0$.  The contact identity implies
$\sum_jc_j=0$.  Small perturbations $p_j\pm tc_j$ again decompose $G$ inside
the fiber.  Conversely, if the evaluation vectors are linearly independent
and $G=tG_1+(1-t)G_2$ with $G_1,G_2$ in the fiber, nonnegativity forces both
measures to be supported on $\supp(G)$.  Their masses solve the same full-rank
linear system, so $G_1=G_2=G$.  This proves the characterization and the
$n$-atom bound.

For measurability, let $E_k$ be the weights admitting an extreme optimizer
with exactly $k$ atoms.  For integers $j\ge2$, consider the compact set of
$(w,\theta_1,\ldots,\theta_k,p)$ in
$[1/j,j]^n\times\Theta^k\times\Delta_k$ satisfying
\[
 p_\ell\ge1/j,\qquad
 \sum_{\ell=1}^kp_\ell A(\theta_\ell)=\vhat(w),
 \qquad \det(V^\top V)\ge1/j,
\]
where $V$ has columns $A(\theta_\ell)$.  Its projection onto $w$ is compact.
Taking the countable union over $j$ gives $E_k$.  Hence
$\{w:s_{\mathrm{ext}}(w)\ge k\}=\bigcup_{\ell=k}^nE_\ell$ is Borel for
each $k$.
\end{proof}

\subsection{Support directions and radial localization}

\begin{lemma}[Matrix sensitivity]
\label{lem:matrix-sensitivity}
At almost every $w>0$, the Hessian $H(w)=\nabla^2\Psi(w)$ exists.  At each
such point and for every extreme optimizer with $k$ atoms, there is an
orthogonal projection $Q$ of rank $k-1$ such that
\[
 \diag(w)^{1/2}H(w)\diag(w)^{1/2}\succeq Q.
\]
Consequently \eqref{eq:hessian-summary} holds with rank
$s_{\mathrm{ext}}(w)-1$.
\end{lemma}

\begin{proof}
By \Cref{lem:regularity}, $z=\nabla\Psi$ is locally Lipschitz, so Rademacher's
theorem gives differentiability almost everywhere.  At those points its
derivative is the symmetric positive semidefinite Hessian of the convex
function $\Psi$.

Fix such a point and an extreme optimizer
$G=\sum_{j=1}^kp_j\delta_{\theta_j}$.  Write
$A_j=A(\theta_j)$ and $\widehat v=\vhat(w)$.  Varying the masses while
preserving their sum generates the relative fitted-vector subspace
\[
 \mathcal S=\left\{
 \left(\frac{\sum_jc_jA_{ij}}{\widehat v_i}\right)_{i=1}^n:
 \sum_jc_j=0\right\}.
\]
The evaluation vectors are linearly, hence affinely, independent by
\Cref{lem:extreme}; therefore $\dim(\mathcal S)=k-1$.  Every sufficiently
small element of $\mathcal S$ is feasible in both signs by changing only the
positive masses.  First-order optimality gives $w^\top u=0$ for
$u\in\mathcal S$.

Let $Q$ be the orthogonal projection onto
$\diag(w)^{1/2}\mathcal S$.  For small $h$, choose the feasible relative
variation
$u=\diag(w)^{-1/2}Q\diag(w)^{-1/2}h$.  Expanding the candidate objective at
weight $w+h$ gives
\[
 \Psi(w+h)\ge\Psi(w)+h^\top z(w)
 +h^\top u-\frac12u^\top\diag(w)u+o(\norm{h}_2^2).
\]
Because $w^\top u=0$, the last two displayed quadratic terms reduce to
\[
 \frac12h^\top\diag(w)^{-1/2}Q\diag(w)^{-1/2}h.
\]
Comparison with the second-order expansion of $\Psi$ yields
$H(w)\succeq\diag(w)^{-1/2}Q\diag(w)^{-1/2}$.  Conjugation proves the claim.
The exceptional set is the differentiability set of $z$ and is independent
of the optimizer.  Choose an extreme optimizer attaining
$s_{\mathrm{ext}}(w)$.
\end{proof}

For the localization argument, retain \eqref{eq:relative-set}--\eqref{eq:width}
and write
\[
 \xi=w-\one,\quad A_0=\norm{P\xi}_2,\quad B_0=\norm{\xi}_2,
 \quad \rho=\frac1{16}.
\]
For $u\in\Ucal$, let $t(u)=\log(\one+u)$ and
$F_w(u)=w^\top t(u)$.  For the exact fit, write
$\widehat u(w)=\vhat(w)/v_0-\one$, $t(w)=t(\widehat u(w))$, and
$Z(w)=\xi^\top t(w)$.

\begin{lemma}[Radial localization]
\label{lem:radial}
Let
\begin{equation}\label{eq:localization-event}
 E=\left\{w>0:
\norm{w-\one}_\infty\le\frac1{16},\quad
 A_0\le\frac{\rho}{24},\quad \eta B_0\le\frac{\rho^2}{24}\right\}.
\end{equation}
For $w\in E$, every $0\le\tau<\rho^2/12$ and every $\tau$-optimal fit with
relative vector $u=\vtil/v_0-\one$ satisfy $\norm{u}_2<\rho$ and
\begin{align}
 \norm{u}_2^2&\le36A_0^2+12\eta B_0+12\tau,\label{eq:radial-u}\\
 0\le-\one^\top\log(\one+u)
 &\le37A_0^2+13\eta B_0+13\tau,\label{eq:radial-loss}\\
 \norm{\log(\one+u)}_2^2
 &\le72A_0^2+24\eta B_0+24\tau.\label{eq:radial-log}
\end{align}
For the exact fit, $0\le Z(w)\le37A_0^2+13\eta B_0$.  The set $E$ is
compact and $\norm{t(w)}_\infty<1/8$ on $E$.
\end{lemma}

\begin{proof}
Equal-weight optimality at $v_0$ implies
$\one^\top u\le0$ for every $u\in\Ucal$.  If $\norm{u}_2\le1/4$, then
$\log(1+s)\le s-s^2/3$ coordinatewise.  On $E$, every $w_i\ge15/16$, and
therefore
\[
 F_w(u)\le \xi^\top u-\frac16\norm{u}_2^2
 \le A_0\norm{u}_2+\eta B_0-\frac16\norm{u}_2^2.
\]
Every feasible $u$ on the sphere $\norm{u}_2=\rho$ consequently satisfies
$F_w(u)\le-\rho^2/12$.  A $\tau$-optimal fit has
$F_w(u)\ge-\tau$ because zero is feasible.  If $\norm{u}_2\ge\rho$, convexity
of $\Ucal$ and concavity of $F_w$ allow radial contraction to the sphere,
contradicting $\tau<\rho^2/12$.  Hence $\norm{u}_2<\rho$.

Let $x=\norm{u}_2$.  The preceding objective inequality gives
$x^2\le6(A_0x+\eta B_0+\tau)$.  The bound
$6A_0x\le x^2/2+18A_0^2$ proves \eqref{eq:radial-u}.  For
$\abs{u_i}<\rho$,
\[
 \abs{\log(1+u_i)}\le\frac{\abs{u_i}}{1-\rho},
 \qquad
 \abs{\log(1+u_i)-u_i}\le\frac{u_i^2}{2(1-\rho)}\le u_i^2.
\]
This proves \eqref{eq:radial-log}.  Moreover,
\[
 \xi^\top\log(\one+u)
 \le A_0x+\eta B_0+\frac1{16}x^2
 \le37A_0^2+13\eta B_0+12\tau,
\]
where constants are enlarged after applying Young's inequality and
\eqref{eq:radial-u}.  Since
$F_w(u)=\one^\top\log(\one+u)+\xi^\top\log(\one+u)\ge-\tau$, and equal-weight
optimality gives $-\one^\top\log(\one+u)\ge0$, we obtain
\eqref{eq:radial-loss}.  For the exact fit, $F_w(\widehat u)\ge0$ gives
$Z(w)\ge-\one^\top t(w)\ge0$ and the stated upper bound.

The set $E$ is closed inside $[15/16,17/16]^n$, hence compact.  For the exact
fit, $\norm{t(w)}_\infty\le\norm{\widehat u(w)}_2/(1-\rho)<1/15<1/8$.
\end{proof}

\subsection{The Gamma change of variables}

\begin{lemma}[Determinant moment]
\label{lem:determinant-moment}
Let $E$ be the event in \eqref{eq:localization-event}.  For independent
$W_i\sim\mathrm{Gamma}(\alpha,\alpha)$,
\begin{equation}\label{eq:determinant-moment}
 \EE\left[\bm 1_{\{W\in E\}}
 \left(\frac{17}{9}\right)^{s_{\mathrm{ext}}(W)-1}
 \exp\{-\alpha[Z(W)+\norm{t(W)}_2^2]\}\right]\le1.
\end{equation}
\end{lemma}

\begin{proof}
On the open set
$U=\{w>0:1+t_i(w)>0\ \text{for all }i\}$, define the map in
\eqref{eq:change-map}.  It is locally Lipschitz by \Cref{lem:regularity}.
Suppose $\mathcal F(w)=\mathcal F(w')=y$.  Then
$w_i=y_i/\{1+t_i(w)\}$ and
\[
 (w-w')^\top\{z(w)-z(w')\}
 =-\sum_i\frac{y_i\{t_i(w)-t_i(w')\}^2}
 {[1+t_i(w)][1+t_i(w')]}\le0.
\]
The gradient of the convex function $\Psi$ is monotone, so the left side is
nonnegative.  Hence $t(w)=t(w')$, and the displayed formula gives $w=w'$.
Thus $\mathcal F$ is injective on $U$.

At a Hessian point, let
\[
 A_t=\diag(\one+t(w)),\qquad
 R_w=\diag(w)^{1/2}\nabla^2\Psi(w)\diag(w)^{1/2}.
\]
The derivative $D\mathcal F=A_t+\diag(w)\nabla^2\Psi(w)$ is similar to
$A_t+R_w$.  By \Cref{lem:matrix-sensitivity},
$R_w\succeq P_*$ for an orthogonal projection of rank
$s_{\mathrm{ext}}(w)-1$.  On $E$, \Cref{lem:radial} gives
$7I/8\preceq A_t\preceq9I/8$.  If $P_*=VV^\top$ with orthonormal columns,
then $V^\top A_t^{-1}V\succeq8I/9$.  Positive-definite determinant
monotonicity and the identity between the nonzero eigenvalues of
$A_t^{-1/2}VV^\top A_t^{-1/2}$ and $V^\top A_t^{-1}V$ yield
\begin{equation}\label{eq:jacobian-lower}
 \det D\mathcal F(w)
 \ge\prod_i\{1+t_i(w)\}
 \left(\frac{17}{9}\right)^{s_{\mathrm{ext}}(w)-1}>0.
\end{equation}

Let $p_\alpha$ be the product Gamma density.  Since
$\mathcal F_i(w)=w_i\{1+t_i(w)\}$,
\[
 \frac{p_\alpha(\mathcal F(w))}{p_\alpha(w)}
 =\prod_i\{1+t_i(w)\}^{\alpha-1}
   \exp\left\{-\alpha\sum_iw_it_i(w)\right\}.
\]
Multiplying by \eqref{eq:jacobian-lower} and using
$\log(1+t)\ge t-t^2$ for $\abs{t}\le1/8$ gives, almost everywhere on $E$,
\[
 \frac{p_\alpha(\mathcal F(w))\det D\mathcal F(w)}{p_\alpha(w)}
 \ge\left(\frac{17}{9}\right)^{s_{\mathrm{ext}}(w)-1}
 e^{-\alpha\{Z(w)+\norm{t(w)}_2^2\}}.
\]
The set $E$ is compact and lies inside $U$.  The injective area formula for Lipschitz maps
\citep[Chapter 3]{evansgariepy2015} therefore gives
\[
 \int_Ep_\alpha(\mathcal F(w))\abs{\det D\mathcal F(w)}\,dw
 =\int_{\mathcal F(E)}p_\alpha(y)\,dy\le1.
\]
The determinant is positive almost everywhere on $E$.  Integrating the
previous lower bound proves \eqref{eq:determinant-moment}.
\end{proof}

\subsection{Proof of the effective-dimension theorem}

\begin{proof}[Proof of \Cref{thm:effective}]
Let $\xi=W-\one$ and put $N_q=r\log5+q$.  The Gamma moment-generating
function implies that, for every deterministic unit vector $v$ and
$\abs{\lambda}\le\sqrt\alpha/2$,
\[
 \log\EE e^{\lambda\sqrt\alpha\,v^\top\xi}
 =\alpha\sum_i\left[-\log\left(1-\frac{\lambda v_i}{\sqrt\alpha}\right)
 -\frac{\lambda v_i}{\sqrt\alpha}\right]\le\lambda^2.
\]
Consequently,
$\PP\{\sqrt\alpha\,v^\top\xi>u\}\le e^{-u^2/4}$ for
$0\le u\le\sqrt\alpha$.  A $1/2$-net of the unit sphere in an
$r$-dimensional subspace has at most $5^r$ points and satisfies
$\norm{y}_2\le2\max_v v^\top y$.  Hence
\begin{equation}\label{eq:proj-concentration}
 \PP\{\alpha\norm{P\xi}_2^2>16N_q\}\le e^{-q}
\end{equation}
provided $\alpha\ge4N_q$.

The scalar Gamma Chernoff bound gives
\[
 \PP\{\norm{\xi}_\infty>2\sqrt{Q/\alpha}\}
 \le2n e^{-Q}=e^{-q}
\]
when $\alpha\ge4Q$.  Let $E_q$ be the intersection of these two events.
Then $\PP(E_q^c)\le2e^{-q}$ and, on $E_q$,
\[
 A_0\le4\sqrt{N_q/\alpha},\qquad
 B_0\le2\sqrt{nQ/\alpha},\qquad
 \eta B_0\le2/\alpha,
\]
where the last inequality uses \eqref{eq:effective-conditions}.  Taking the
universal constant in \eqref{eq:effective-conditions} sufficiently large
ensures $E_q\subset E$ in \eqref{eq:localization-event}.

Apply \Cref{lem:radial}.  Since $N_q\le C(r+q)$, its exact-fit bounds yield
\[
 0\le\one^\top\log v_0-\one^\top\log\vhat(W)
 \le C\frac{r+q}{\alpha},
 \qquad
 \norm{\log\{\vhat(W)/v_0\}}_2^2
 \le C\frac{r+q}{\alpha}.
\]
The same lemma gives \eqref{eq:effective-approx-fit} simultaneously for all
$\tau<c$, after taking $c=\rho^2/12$ and enlarging $C$.

For the support bound, \Cref{lem:radial} also implies on $E_q$
\[
 \alpha\{Z(W)+\norm{t(W)}_2^2\}
 \le109\alpha A_0^2+37\alpha\eta B_0
 \le1744N_q+74=:H_q.
\]
Set $\kappa=\log(17/9)$.  By \Cref{lem:determinant-moment} and Markov's
inequality,
\[
 \PP\{E_q,\ (s_{\mathrm{ext}}(W)-1)\kappa>H_q+q\}\le e^{-q}.
\]
Adding $\PP(E_q^c)$ shows that, outside probability $3e^{-q}$,
\[
 s_{\mathrm{ext}}(W)
 \le1+\frac{1744N_q+74+q}{\kappa}
 \le C(r+q).
\]
All conclusions hold on the same event.
\end{proof}

\section{Gaussian approximation and uniqueness}
\label{app:gaussian}

\subsection{Proof of the Gaussian relative-width bound}

\begin{proof}[Proof of \Cref{prop:gaussian-width}]
Set $b_i=\phi_d(x_i)/v_{0,i}$.  The Gaussian identity gives
\begin{equation}\label{eq:gaussian-ratio}
 \frac{\phi_d(x_i-\theta)}{v_{0,i}}
 =b_i e^{-\norm{\theta}^2/2}e^{x_i^\top\theta}.
\end{equation}
Because
\[
 \frac{v_{0,i}}{\phi_d(x_i)}
 =\int e^{x_i^\top\vartheta-\norm{\vartheta}^2/2}\,dG_0(\vartheta)
 \ge e^{-ST-S^2/2},
\]
we have $b_i\le e^{ST+S^2/2}$.  For $\abs{t}\le ST$, Taylor's theorem yields
\[
 \left|e^t-\sum_{j=0}^m\frac{t^j}{j!}\right|
 \le e^{ST}\frac{(ST)^{m+1}}{(m+1)!}.
\]
Multiplying by the remaining factors in \eqref{eq:gaussian-ratio} shows that
\[
 b_i e^{-\norm{\theta}^2/2}
 \sum_{j=0}^m\frac{(x_i^\top\theta)^j}{j!}
\]
approximates $\phi_d(x_i-\theta)/v_{0,i}$ with coordinatewise error at most
$\eps_m$ in \eqref{eq:eps-m}, uniformly over $i$ and $\theta\in K$.

The multinomial expansion of $(x_i^\top\theta)^j$ uses only monomials
$x_i^\beta$ with multi-index $\abs{\beta}\le m$.  Hence every truncated vector
lies in the span of
\[
 \left\{(b_ix_i^\beta)_{i=1}^n:\abs{\beta}\le m\right\},
\]
whose dimension is at most $\binom{m+d}{d}$.  Adjoin $\one$ to this span and
let $P_m$ be the orthogonal projection onto the resulting subspace.  Its rank
obeys \eqref{eq:rank-m}.  Integrating the coordinatewise approximation under
any $G\in\Pcal(K)$ preserves its error, and orthogonal projection minimizes
Euclidean distance to the subspace.  This proves \eqref{eq:eta-m}.  If
$ST=0$, the exponential term is constant in the relevant inner product and
the expansion is exact.
\end{proof}

\subsection{Proof of the conditional Gaussian corollary}

\begin{proof}[Proof of \Cref{cor:conditional}]
Choose any ordinary NPMLE $\Ghat_0$ and let
$v_0=(f_{\Ghat_0}(x_i))_{i=1}^n$.  This vector is independent of the chosen
maximizing representation by \Cref{lem:regularity}.  Apply
\Cref{prop:gaussian-width} with
\[
 m=\left\lceil C_m\frac{\log n}{\log\log n}\right\rceil,
\]
where $C_m>8$ is a sufficiently large constant depending only on the fixed
parameters.  The factorial inequality $(m+1)!\ge\{(m+1)/e\}^{m+1}$ and
$T\le C_0\sqrt{\log n}$ give
\begin{align*}
 \log\eps_m
 &\le 2ST+S^2/2+(m+1)\log\frac{eST}{m+1}\\
 &\le-3\log n
\end{align*}
for all sufficiently large $n$.  Indeed,
$\log(ST)\le\tfrac12\log\log n+O(1)$, whereas
$\log(m+1)=\log\log n-\log\log\log n+O(1)$.  If $ST=0$, the error is zero.
Thus
\[
 r:=\rank(P_m)=O(r_n),\qquad
 \eta:=\eta_m\le n^{-5/2}.
\]

Take $q=(b+2)\log n$ and $Q=q+\log(2n)$.  Then
$r+Q=O_b(\bar r_n)$ and
\[
 \frac{\alpha_n}{r+Q}\asymp_b\log n\to\infty.
\]
Moreover,
\[
 \eta\sqrt{\alpha_n nQ}
 \le n^{-5/2}\sqrt{\bar r_n\log n\cdot nQ}=o(1).
\]
The hypotheses of \Cref{thm:effective} therefore hold.  Its failure
probability is at most $3n^{-b-2}$, and it yields
$s_{\mathrm{ext}}(W)\le C_b\bar r_n$ together with the likelihood and
fitted-vector conclusions.

For the coordinate weights, the Gamma Chernoff inequality
\begin{equation}\label{eq:gamma-coordinate}
 \PP\{\abs{W_i-1}>t\}\le2e^{-\alpha_nt^2/4},
 \qquad 0<t\le1,
\end{equation}
follows from the Gamma moment-generating function.  Choose
$t=C_b'/\sqrt{\bar r_n}$ with $(C_b')^2/4\ge b+3$.  Since
$\alpha_nt^2=(C_b')^2\log n$, a union bound gives
\[
 \PP\{\max_i\abs{W_i-1}>t\}
 \le2n^{1-(C_b')^2/4}\le2n^{-b-2}.
\]
Intersect the events and enlarge constants.
\end{proof}

\subsection{Proof of generic evaluation independence}
\label{app:generic}

\begin{proof}[Proof of \Cref{prop:generic}]
Fix an integer $k$ satisfying $k(d+1)\le n$.  A nonzero linear dependence
among $k$ Gaussian evaluation columns has a coefficient vector with at least
one nonzero coordinate.  Normalize that coordinate to one; there are only
$k$ possible choices.  In one such chart, collect the $kd$ location
coordinates and the remaining $k-1$ coefficients in a parameter
$\eta\in U\subset\RR^p$, where
\[
 p=kd+k-1=k(d+1)-1<n.
\]
The domain $U$ is open after imposing distinctness of the locations.  For the
normalized coefficient vector $(c_1,\ldots,c_k)$, define
\[
 F_\eta(y)=\sum_{j=1}^kc_j\phi_d(y-\theta_j).
\]
This function is jointly real analytic in $(\eta,y)$.

For every admissible $\eta$, $F_\eta$ is not identically zero.  Choose
$v\in\RR^d$ such that the scalars $v^\top\theta_j$ are distinct; the excluded
$v$ lie in finitely many proper hyperplanes.  If $F_\eta(tv)=0$ for every
$t$, then after division by the positive common Gaussian factor,
\[
 \sum_{j=1}^kc_j e^{-\norm{\theta_j}^2/2}e^{t v^\top\theta_j}=0
 \qquad\text{for all }t.
\]
Differentiating at zero for orders $0,\ldots,k-1$ gives a Vandermonde system
with distinct nodes $v^\top\theta_j$, forcing every coefficient to vanish.
This contradicts the normalization.

Consider the incidence set
\[
 I=\{(\eta,x_1,\ldots,x_n)\in U\times(\RR^d)^n:
        F_\eta(x_i)=0\text{ for all }i\}.
\]
At an incidence point, analyticity and nontriviality imply that for each $i$
there is a multi-index $\alpha_i$ of minimal total order with
$\partial_y^{\alpha_i}F_\eta(x_i)\ne0$.  Its order is at least one because
$F_\eta(x_i)=0$.  Choose a coordinate $\ell_i$ with
$\alpha_{i,\ell_i}>0$ and put $\beta_i=\alpha_i-e_{\ell_i}$.  By minimality,
\[
 H_i(\eta,x_i):=\partial_y^{\beta_i}F_\eta(x_i)=0,
 \qquad
 \partial_{x_{i,\ell_i}}H_i(\eta,x_i)\ne0.
\]
For each fixed choice of the countably many derivative indices, restrict to
the open region on which the displayed derivatives are nonzero.  The
Jacobian of $(H_1,\ldots,H_n)$ with respect to the selected data coordinates
$(x_{1,\ell_1},\ldots,x_{n,\ell_n})$ is diagonal and nonsingular.  The
implicit-function theorem therefore places the corresponding part of $I$ in
a smooth submanifold of codimension $n$ in
$U\times\RR^{nd}$.  Its dimension is
\[
 nd+p-n<nd.
\]
Each such submanifold has a countable atlas.  On compact subsets of a chart,
projection to the $nd$ data coordinates is Lipschitz; the image of a
$q$-dimensional bounded set under a Lipschitz map has zero
$nd$-dimensional Lebesgue measure whenever $q<nd$.  Thus the projection of
$I$ is contained in a countable union of null sets.

Taking the finite union over coefficient normalizations and over admissible
$k$ produces a Borel null set $N_{n,d}$.  Let
$\mathcal X_{n,d}=(\RR^d)^n\setminus N_{n,d}$.  This proves the simultaneous
column-independence assertion.

For the final statement, suppose every extreme optimizer has at most $m$
atoms and $2m(d+1)\le n$.  If two distinct extreme optimizers existed, their
difference would be a nonzero signed measure on at most $2m$ distinct
locations and would have zero Gaussian evaluation vector.  After deleting
zero coefficients, this would contradict the established independence.
Hence the compact convex optimizer fiber has at most one extreme point.  By
the Krein--Milman theorem it is the closed convex hull of its nonempty set of
extreme points, so it must be a singleton.  Its unique member is extreme and
has at most $m$ atoms.
\end{proof}

\section{Likelihood nets and Hellinger transfer}
\label{app:hellinger}

\subsection{A finite Gaussian-mixture likelihood net}

\begin{lemma}[Simultaneous Hellinger and log-likelihood net]
\label{lem:finite-net}
Fix $d\ge1$, a nonempty compact $K\subset B(0,S)$, and $b>0$.  For all
sufficiently large $n$, there are a deterministic finite set
$\Ncal_n\subset\{f_G:G\in\Pcal(K)\}$ and a deterministic
$T_n=O(\sqrt{\log n})$ such that
\begin{equation}\label{eq:net-size}
 \log\abs{\Ncal_n}
 \le C\frac{(\log n)^{d+1}}{(\log\log n)^d},
\end{equation}
\begin{equation}\label{eq:radius-prob}
 \sup_{G^*\in\Pcal(K)}
 \PP_{G^*}\left\{\max_{i\le n}\norm{X_i}>T_n\right\}
 \le n^{-b-2},
\end{equation}
and every $G\in\Pcal(K)$ admits $p\in\Ncal_n$ satisfying
\begin{equation}\label{eq:net-approx}
 H(p,f_G)\le n^{-3},
 \qquad
 \sup_{\norm{x}\le T_n}\abs{\log p(x)-\log f_G(x)}\le n^{-3}.
\end{equation}
Constants are uniform over $G$ and $G^*$.
\end{lemma}

\begin{proof}
For $n\ge3$, take
\[
 T_n=S+\sqrt{2d\log(2d n^{b+3})}.
\]
Under any $G^*\in\Pcal(K)$, $X_i=\Theta_i+Z_i$ with
$\norm{\Theta_i}\le S$.  The standard Gaussian coordinate bound and a union
bound over $nd$ coordinates prove \eqref{eq:radius-prob}.  If $S=0$, the class
contains only $\phi_d$, so assume $S>0$.

Let
\[
 m=\left\lceil C_0\frac{\log n}{\log\log n}\right\rceil,
 \qquad k=\binom{2m+d}{d},
\]
with fixed $C_0>16$.  Consider the vector of all nonconstant monomials in
$\theta$ of total degree at most $2m$.  Its image of $K$ is compact, and the
vector of its integrals under any $G\in\Pcal(K)$ lies in the corresponding
convex hull.  Carath\'eodory's theorem gives a measure $G_m$ with at most $k$
atoms that matches every moment of total degree at most $2m$.

For $\norm{x}\le T_n$, set $L=ST_n+S^2/2$.  The exponent
$x^\top\theta-\norm{\theta}^2/2$ has absolute value at most $L$ on $K$, and
$f_G(x)/\phi_d(x)\ge e^{-L}$.  Its Taylor polynomial through order $m$ has
total degree at most $2m$ in $\theta$, so its integrals under $G$ and $G_m$
coincide.  Bounding both remainders yields
\begin{equation}\label{eq:relative-approx}
 \sup_{\norm{x}\le T_n}
 \left|\frac{f_{G_m}(x)}{f_G(x)}-1\right|
 \le2e^{2L}\frac{L^{m+1}}{(m+1)!}\le n^{-8}
\end{equation}
for all large $n$.  The last inequality follows because
$L=O(\sqrt{\log n})$ and
$(m+1)!\ge\{(m+1)/e\}^{m+1}$.

Moment matching also gives a global approximation.  Let
$\Delta=G-G_m$.  Completing the square shows
\[
 \int_{\RR^d}
 \frac{\phi_d(x-\theta)\phi_d(x-\eta)}{\phi_d(x)}\,dx
 =e^{\theta^\top\eta}.
\]
The exponential series is uniformly convergent on $K\times K$.  Every term
of degree at most $2m$ integrates to zero against $\Delta\otimes\Delta$, and
$\norm{\Delta}_{\TV}\le2$.  Therefore
\begin{align}
 Q_m:=\int\frac{(f_G-f_{G_m})^2}{\phi_d}
 &=\iint e^{\theta^\top\eta}\,d\Delta(\theta)d\Delta(\eta)\notag\\
 &\le4\sum_{j>2m}\frac{S^{2j}}{j!}
 \le4e^{S^2}\frac{S^{4m+2}}{(2m+1)!}\le n^{-32}.
 \label{eq:Qm}
\end{align}
By Cauchy--Schwarz,
$H^2(f_G,f_{G_m})\le\norm{f_G-f_{G_m}}_1\le Q_m^{1/2}$, so
$H(f_G,f_{G_m})\le n^{-8}$.

It remains to discretize the locations and masses.  Put $\eps=n^{-16}$.
Partition a fixed bounding cube into cells of diameter at most $\eps$ and
choose one point of $K$ from every cell that meets $K$.  The resulting
location net has size $J_n\le C_Kn^{16d}$.  Write
$G_m=\sum_{j=1}^k\lambda_j\delta_{\theta_j}$, padding with zero masses if
needed, and round every location to the net.  Let
$M=\lceil kn^{16}\rceil$.  For $j<k$, put
$\lambda_j'=\lfloor M\lambda_j\rfloor/M$, and let
$\lambda_k'=1-\sum_{j<k}\lambda_j'$.  Then
$\sum_j\abs{\lambda_j'-\lambda_j}\le2\eps$.  Let $G'$ be the resulting
discrete measure and $p=f_{G'}$.

Moving a location by at most $\eps$ changes its Gaussian log kernel on
$\{\norm{x}\le T_n\}$ by at most $(T_n+S)\eps$.  Weight rounding changes the
mixture by relative amount at most $2\eps e^{2L}$.  Combining these estimates
with \eqref{eq:relative-approx}, and using
$\abs{\log(1+u)}\le2\abs{u}$ for $\abs{u}\le1/2$, gives
\[
 \sup_{\norm{x}\le T_n}\abs{\log p(x)-\log f_G(x)}\le n^{-3}
\]
for all large $n$.  Globally, the fundamental theorem of calculus and
$\int\norm{\nabla\phi_d(x)}\,dx<\infty$ give an $L^1$ error $O_d(\eps)$ from
location rounding; weight rounding adds at most $2\eps$.  The Hellinger
triangle inequality therefore gives $H(p,f_G)\le n^{-3}$.

Let $\Ncal_n$ contain the densities generated by every ordered list of $k$
location-net points and every simplex-grid mass vector with denominator $M$.
Then each $G$ has a point satisfying \eqref{eq:net-approx}, and
\[
 \log\abs{\Ncal_n}
 \le k\{\log J_n+\log(M+1)\}
 =O(k\log n)
 =O\left\{\frac{(\log n)^{d+1}}{(\log\log n)^d}\right\}.
\]
This proves \eqref{eq:net-size}.
\end{proof}

\subsection{Proof of the near-MLE theorem}

\begin{proof}[Proof of \Cref{thm:near-mle}]
Fix $G^*\in\Pcal(K)$ and write $p_*=f_{G^*}$.  Let $\Ncal_n,T_n$ be from
\Cref{lem:finite-net}, and let
$E_n=\{\max_i\norm{X_i}\le T_n\}$.  On $E_n$, for any candidate satisfying
\eqref{eq:near-truth}, choose $p\in\Ncal_n$ as in
\eqref{eq:net-approx}.  Then
\begin{equation}\label{eq:net-lr-lower}
 \sum_{i=1}^n\log\frac{p(X_i)}{p_*(X_i)}
 \ge-1-n^{-2}\ge-2.
\end{equation}

For deterministic $p\in\Ncal_n$, independence gives
\begin{align*}
 \EE_{p_*}\exp\left\{\frac12\sum_{i=1}^n
       \log\frac{p(X_i)}{p_*(X_i)}\right\}
 &=\left(\int\sqrt{pp_*}\right)^n\\
 &=\left\{1-\frac12H^2(p,p_*)\right\}^n
 \le e^{-nH^2(p,p_*)/2}.
\end{align*}
Markov's inequality therefore yields
\begin{equation}\label{eq:fixed-net-test}
 \PP_{p_*}\left\{\sum_i\log\frac{p(X_i)}{p_*(X_i)}\ge-2\right\}
 \le\exp\left\{1-\frac n2H^2(p,p_*)\right\}.
\end{equation}
Set
\[
 A_n=\log\abs{\Ncal_n}+(b+2)\log n+2,
 \qquad t_n=\sqrt{2A_n/n}.
\]
Summing \eqref{eq:fixed-net-test} over all net points with
$H(p,p_*)\ge t_n$ shows that, outside an event of probability at most
$e^{-1}n^{-b-2}$, every such point has log-likelihood ratio below $-2$.

Intersect this event with $E_n$.  The net point associated with any candidate
in \eqref{eq:near-truth} satisfies \eqref{eq:net-lr-lower}, so it must obey
$H(p,p_*)<t_n$.  Hence, simultaneously over all candidates,
\[
 H^2(f_G,p_*)\le(t_n+n^{-3})^2
 \le\frac{4A_n}{n}+2n^{-6}.
\]
The entropy bound \eqref{eq:net-size} gives the claimed rate.  The radius and
likelihood-test exceptional probabilities are uniform over
$G^*\in\Pcal(K)$, completing the proof.
\end{proof}

\section{Dirichlet representation and auxiliary concentration}
\label{app:dirichlet}

\begin{proof}[Proof of \Cref{prop:dirichlet}]
The standard Gamma--Dirichlet change of variables gives
\[
 (P_1,\ldots,P_n)\sim\mathrm{Dirichlet}(\alpha_n,\ldots,\alpha_n),
 \qquad T_W\sim\mathrm{Gamma}(n\alpha_n,\alpha_n),
\]
and independence between $P$ and $T_W$.  The Dirichlet moments are
\[
 \EE P_i=\frac1n,
 \quad
 \operatorname{Var}(P_i)=\frac{n-1}{n^2(n\alpha_n+1)},
 \quad
 \operatorname{Cov}(P_i,P_j)=-\frac1{n^2(n\alpha_n+1)}\quad(i\ne j).
\]
Substitution gives the conditional mean and variance of
$F_n^{(\alpha_n)}h=\sum_iP_ih(X_i)$.  Also,
\[
 \norm{F_n^{(\alpha_n)}-F_n}_{\TV}
 \le\frac12\sum_i\abs{P_i-1/n}.
\]
Cauchy--Schwarz and the marginal Dirichlet variance therefore imply
\[
 \EE\bigl(\norm{F_n^{(\alpha_n)}-F_n}_{\TV}\mid X\bigr)
 \le\frac n2\sqrt{\operatorname{Var}(P_1)}
 =\frac12\sqrt{\frac{n-1}{n\alpha_n+1}}.
\]

Because $\EE W_i=1$ and $\EE(P_i\mid X)=1/n$, independence and
integrability give
\[
 \EE_{X,W}\frac1n\sum_iW_i\log f_G(X_i)
 =\EE_{G^*}\log f_G(X_1)
 =\EE_{X,W}\sum_iP_i\log f_G(X_i).
\]
For compact $K\subset B(0,S)$,
\[
 e^{-S\norm{x}-S^2/2}
 \le\frac{f_G(x)}{\phi_d(x)}\le e^{S\norm{x}},
\]
so all displayed log densities are integrable under a compactly supported
Gaussian mixture.  Their population difference at $G^*$ and $G$ is
\[
 \int f_{G^*}(x)\log\frac{f_{G^*}(x)}{f_G(x)}\,dx
 =\KL(f_{G^*}\|f_G).
\]
\end{proof}

\begin{lemma}[Gamma coordinate concentration]
\label{lem:gamma-concentration}
If $W\sim\mathrm{Gamma}(a,a)$, then for $0<t\le1$,
\[
 \PP\{\abs{W-1}>t\}\le2e^{-at^2/4}.
\]
\end{lemma}

\begin{proof}
For $0<\lambda<a$,
$\EE e^{\lambda(W-1)}=e^{-\lambda}(1-\lambda/a)^{-a}$.  Optimizing the
Chernoff bound gives
$\PP(W\ge1+t)\le\exp\{-a(t-\log(1+t))\}$.  Since
$t-\log(1+t)\ge t^2/4$ for $0\le t\le1$, this is at most $e^{-at^2/4}$.
Similarly,
$\PP(W\le1-t)\le\exp\{-a(-t-\log(1-t))\}\le e^{-at^2/2}$.
Adding the two tails proves the claim.
\end{proof}

\begin{thebibliography}{99}

\bibitem[Evans and Gariepy(2015)]{evansgariepy2015}
Evans, L. C. and Gariepy, R. F. (2015).
\newblock \emph{Measure Theory and Fine Properties of Functions}, revised ed.
\newblock Boca Raton: CRC Press.

\bibitem[Jia et al.(2023)Jia, Polyanskiy and Wu]{jiaetal2023}
Jia, Z., Polyanskiy, Y. and Wu, Y. (2023).
\newblock Entropic characterization of optimal rates for learning Gaussian
mixtures.
\newblock In \emph{Proceedings of the 36th Conference on Learning Theory},
PMLR 195, 4296--4335.

\bibitem[Kiefer and Wolfowitz(1956)]{kieferwolfowitz1956}
Kiefer, J. and Wolfowitz, J. (1956).
\newblock Consistency of the maximum likelihood estimator in the presence of
infinitely many incidental parameters.
\newblock \emph{Annals of Mathematical Statistics}, \textbf{27}, 887--906.

\bibitem[Kim and Guntuboyina(2022)]{kimguntuboyina2022}
Kim, A. K. H. and Guntuboyina, A. (2022).
\newblock Minimax bounds for estimating multivariate Gaussian location
mixtures.
\newblock \emph{Electronic Journal of Statistics}, \textbf{16}, 1461--1484.

\bibitem[Koenker and Mizera(2014)]{koenkermizera2014}
Koenker, R. and Mizera, I. (2014).
\newblock Convex optimization, shape constraints, compound decisions, and
empirical Bayes rules.
\newblock \emph{Journal of the American Statistical Association},
\textbf{109}, 674--685.

\bibitem[Laird(1978)]{laird1978}
Laird, N. (1978).
\newblock Nonparametric maximum likelihood estimation of a mixing
distribution.
\newblock \emph{Journal of the American Statistical Association},
\textbf{73}, 805--811.

\bibitem[Lindsay(1983)]{lindsay1983geometry}
Lindsay, B. G. (1983).
\newblock The geometry of mixture likelihoods: a general theory.
\newblock \emph{Annals of Statistics}, \textbf{11}, 86--94.

\bibitem[Newton and Raftery(1994)]{newtonraftery1994}
Newton, M. A. and Raftery, A. E. (1994).
\newblock Approximate Bayesian inference with the weighted likelihood
bootstrap.
\newblock \emph{Journal of the Royal Statistical Society, Series B},
\textbf{56}, 3--26.

\bibitem[Polyanskiy and Wu(2020)]{polyanskiywu2020}
Polyanskiy, Y. and Wu, Y. (2020).
\newblock Self-regularizing property of nonparametric maximum likelihood
estimator in mixture models.
\newblock arXiv:2008.08244.

\bibitem[Rubin(1981)]{rubin1981}
Rubin, D. B. (1981).
\newblock The Bayesian bootstrap.
\newblock \emph{Annals of Statistics}, \textbf{9}, 130--134.

\bibitem[Saha and Guntuboyina(2020)]{sahaguntuboyina2020}
Saha, S. and Guntuboyina, A. (2020).
\newblock On the nonparametric maximum likelihood estimator for Gaussian
location mixture densities with application to Gaussian denoising.
\newblock \emph{Annals of Statistics}, \textbf{48}, 738--762.

\bibitem[Soloff et al.(2025)Soloff, Guntuboyina and Sen]{soloffetal2025}
Soloff, J. A., Guntuboyina, A. and Sen, B. (2025).
\newblock Multivariate, heteroscedastic empirical Bayes via nonparametric
maximum likelihood.
\newblock \emph{Journal of the Royal Statistical Society, Series B},
\textbf{87}, 1--32.

\bibitem[Wang(2026)]{wang2026}
Wang, H. (2026).
\newblock On finite Gaussian mixtures: finiteness of the number of modes and
an application to NPMLE.
\newblock arXiv:2608.16675.

\bibitem[Zhang(2009)]{zhang2009}
Zhang, C.-H. (2009).
\newblock Generalized maximum likelihood estimation of normal mixture
densities.
\newblock \emph{Statistica Sinica}, \textbf{19}, 1297--1318.

\end{thebibliography}

\end{document}
